\section{/ 7.3 / 7.4 / 7.6 Analysis}\label{sec-7-2}

\subsection{Scaffolding Pays --- but only on free-text + non-reasoning backbones}\label{scaffolding-pays-but-only-on-free-text-non-reasoning-backbones}

Phase 4a holds the central backbone constant and varies scaffold complexity:
- \textbf{No scaffold} (llm\_control): single LLM call, structured output.
- \textbf{Single-pass multi-agent} (mdagents): 3-domain experts + Chief MO synthesis.
- \textbf{Multi-round debate} (medagents): 3-expert iterative refinement.
- \textbf{OSCE simulation} (agentclinic): doctor / patient / measurement / moderator loop.
- \textbf{Panel orchestration} (MAI-DxO): \textasciitilde5 medical agents with ordering / questioning.

\textbf{On PP-Store with Gemini Flash} (N=500) the scaffolding ladder shows:
- llm\_control: 0.31
- mdagents (intermediate): 0.32 (+1 pp)
- \textbf{medagents (debate): 0.33 (+1 pp over mdagents, +2 pp over llm\_control)}
- agentclinic (OSCE): 0.25 (-6 pp regression vs control)
- maidxo (panel): 0.02 (catastrophic, see \S{} below)

\textbf{Interpretation}: scaffolding that \emph{deepens} deliberation on the same input
(medagents debate) helps only marginally (+2 pp, within overlapping CIs at
N=500). Scaffolding that \emph{changes the input format} (agentclinic OSCE dialogue,
maidxo panel orchestration) regresses on HPO-list input because the agent's
design assumes narrative input. (The much larger +5--7 pp gap reported in the v0
draft was an artifact of stale N=50 medagents numbers.)

\textbf{Backbone interaction}: on DS V4-Pro mdagents reaches 0.35 (best of its
backbones, N=100), suggesting V4-Pro pairs well with mdagents' moderator-vote
architecture. On GPT-5-minimal mdagents drops to 0.28 --- the moderator's ``weigh
expert opinions'' prompt benefits from GPT-5's reasoning, which we forced off.

\textbf{MAI-DxO catastrophic failure (R@1 \ensuremath{\leq} 0.07 across all backbones)}:
MAI-DxO is designed for NEJM clinicopathologic case-reports (\ensuremath{\geq}2000-word
narratives). On HPO-list input the panel's ``ask the patient'' mechanism
degenerates --- the input \emph{is} the answer to most questions. The panel
emits vitals / lab values (DLCO, LVEF, blood pressure) as ranked
candidates, and we apply a 13-pattern noise filter that catches them but
can't compensate. Documented in \Cref{sec-5-1}.

\subsection{Genotype Channel Helps Any Agent that Ingests It (+20 pp)}\label{sec-7-3}

Phase 3.2 P3 (HPO + structured variants), full-N paired on 500 PP-Store
cases with \ensuremath{\geq}1 structured variant (2026-07-06; llm\_control, Gemini Flash, same
cases both modes):

\begin{table*}[t]
\centering\scriptsize\setlength{\tabcolsep}{3.5pt}
\caption{P2\(\rightarrow\)P3 genotype-aware lift, per agent.}\label{tbl:p3lift}
\begin{tabular}{@{}
llll@{}}
\toprule
Agent & P2 (HPO-only) R@1 & P3 (HPO + variants) R@1 & Lift \\
\midrule
llm\_control (n=500 paired) & 0.296 & \textbf{0.494} & \textbf{+19.8 pp} \\
deeprare (n=50 pilot) & 0.22 & \textbf{0.38} & +16 pp \\
\bottomrule
\end{tabular}
\end{table*}

\textbf{H2 confirmed at full-N and FWE-robust} (\Cref{sec-8-8}): the +19.8 pp lift on n=500
paired cases gives a McNemar \ensuremath{\chi}\textsuperscript{2}(cc)=85 (P3-win 106 vs P2-win 7) and 2-prop
z=6.40 (Holm-adj p=3.0e-10) --- up from the earlier n=50 pilot (z=2.08) that did
not survive correction. Both agents gain \textasciitilde20 pp R@1 from a structured-text
variants block in the prompt. The lift is not DeepRare-specific ---
llm\_control absorbs the same lift, suggesting any LLM can leverage variants when
given them in a parseable form.

The 28-pp gap to DeepRare's published HPO+VCF 70.6 \% is explained by
three documented setup differences:
- Our variants are a structured-text block, not real VCF integrated through
DeepRare's Phenotype Tool.
- Web search disabled (\texttt{DEEPRARE\_NO\_WEB=1}) for contamination control.
- Phenopacket-Store is a harder mixed-difficulty corpus.

\textbf{Honest framing}: variant channel adds \textasciitilde20 pp R@1 to any agent that
ingests it, \emph{not} ``DeepRare specifically exploits genotype-aware
reasoning'' --- the latter claim does not survive contact with the LLM
control baseline.

\subsection{Faithfulness vs Accuracy --- Decoupled (H10 supported)}\label{sec-7-4}

Phase 1 P5 reasoning\_communication pilot on 50-case sample (LLM-judge
faithfulness scoring with Gemini-judge then Claude-judge for bias
detection, \Cref{sec-7-5}):

(Pending Phase 4a P5 data --- to be filled when final.)

\textbf{Preliminary finding from Phase 1}: Spearman \ensuremath{\rho} between accuracy R@1 and
faithfulness score is 0.36 \ensuremath{\pm} 0.11 (95\% CI from bootstrap), well below
the 0.5 threshold pre-registered as ``decoupled.'' This is the strongest
single argument that accuracy-only benchmarks under-evaluate rare-disease diagnostic AI --- high-accuracy agents can have low faithfulness scores
(stating high confidence without justification, hallucinating differential
reasoning), and vice versa.

\subsection{Dataset Difficulty Stratification}\label{sec-7-6}

Four-layer dataset selection deliberately spans difficulty:

\begin{table*}[t]
\centering\scriptsize\setlength{\tabcolsep}{3.5pt}
\caption{Best LLM vs.~best classical/offline R@1, per layer (H1).}\label{tbl:h1layer}
\begin{tabular}{@{}
>{\raggedright\arraybackslash}p{(\linewidth - 4\tabcolsep) * \real{0.3333}}
  >{\raggedright\arraybackslash}p{(\linewidth - 4\tabcolsep) * \real{0.3333}}
  >{\raggedright\arraybackslash}p{(\linewidth - 4\tabcolsep) * \real{0.3333}}@{}}
\toprule
\begin{minipage}[b]{\linewidth}\raggedright
Layer
\end{minipage} & \begin{minipage}[b]{\linewidth}\raggedright
Best LLM R@1 (N=500)
\end{minipage} & \begin{minipage}[b]{\linewidth}\raggedright
Best Classical/Offline R@1
\end{minipage} \\
\midrule
Phenopacket-Store (HPO+demographic, curated rare diseases) & 0.33 (medagents Gemini) & \textbf{0.46} (lirical) \\
RareArena RDS (free-text vignette, narrative) & 0.32 (medagents Gemini) & n/a (no HPO) \\
RareBench HF (HPO-only, sparse, expert curated) & 0.12 (mdagents Gemini) & \textbf{0.35} (vc\_rdagent offline) \\
MIMIC-IV diverse (structured note \ensuremath{\rightarrow} named disease) & 0.39 (mdagents Gemini) & n/a \\
MIMIC rd\_detection prompt (pilot reframe) & \textbf{0.56} (extracts named rare disease from list) & n/a \\
\bottomrule
\end{tabular}
\end{table*}

\textbf{Reading}:
- \textbf{PP-Store is the ``easiest'' layer} --- curated cases, expert-cleaned HPO,
paper-faithful baselines (lirical replicates paper-claimed 0.42 \ensuremath{\pm} 0.05).
- \textbf{RareBench HF is hardest for LLM} --- universal \ensuremath{\leq}0.09 R@1 despite agents
emitting reasonable named diagnoses; the ID-mapping cross-ref limits
evaluator credit. Two paths for paper: (a) report strict and add
hierarchy-aware secondary metric, (b) drop RareBench from main table
and use it only as ``ID-disambiguation stress test'' in the appendix.
- \textbf{MIMIC} under default DDx prompt scores low because input ICD codes
conflate primary rare-disease with comorbidities; the rd\_detection
reframe (\Cref{sec-9} L4 follow-up) recovers to 0.56 by changing task to
``identify which condition is the rare disease.''

\subsection{H1 --- Prevalence stratification (real Orphanet prevalence)}\label{sec-7-7}

Pre-registered H1: R@1 declines monotonically from common-rare to super-rare.
Tested with real Orphadata prevalence (5,108 ORPHA
codes; point-prevalence preferred, rarest validated class per disease), gold
mapped via direct ORPHA or OMIM\ensuremath{\rightarrow}ORPHA cross-map. Pooled R@1 by tier
(commonest\ensuremath{\rightarrow}rarest):

\begin{table*}[t]
\centering\scriptsize\setlength{\tabcolsep}{3.5pt}
\caption{Prevalence-tier R@1: LLM vs.~classical (H1).}\label{tbl:h1tier}
\begin{tabular}{@{}
lll@{}}
\toprule
Tier & LLM (Gemini, N=500 cells) & Classical (LIRICAL+VC-RDAgent) \\
\midrule
common-rare (\ensuremath{\geq}1/10,000) & 0.37 (n=156) & 0.30 (n=64) \\
moderate (1-9/100,000) & 0.26 (n=690) & 0.23 (n=347) \\
ultra-rare (1-9/1,000,000) & 0.39 (n=693) & 0.33 (n=322) \\
\textbf{super-rare (\textless1/1,000,000)} & \textbf{0.22} (n=1167) & \textbf{0.50} (n=529) \\
\bottomrule
\end{tabular}
\end{table*}

\textbf{Strict monotonic H1 is \emph{not} supported} for either class (an ultra-rare
mid-spike breaks monotonicity). But the tail contrast is the real story:
- \textbf{LLMs decline toward the rarest tier} (common 0.37 \ensuremath{\rightarrow} super-rare 0.22,
\ensuremath{-}15 pp) --- consistent with the training-frequency-exposure mechanism H1
posits, just non-monotonic in the middle.
- \textbf{Classical/offline agents do their \emph{best} on super-rare disease} (0.50,
their top tier) --- the inverse of H1. LIRICAL's Bayesian likelihood and
VC-RDAgent's information-content weighting reward the highly specific
phenotype fingerprints that ultra-rare diseases present.
- On the rarest tier the classical-vs-LLM gap widens to +28 pp (0.50 vs
0.22), strengthening F1: the rarer the disease, the larger the classical
advantage.

\textbf{Operationalization note (for PI review)}: ``rarest validated class per
disease'' is a conservative choice when a disorder has multiple prevalence
estimates; point-prevalence entries are preferred over cases/families.
This supersedes the sample-frequency proxy in A10 (which left PP-Store empty
because its golds are OMIM-keyed). Visualised in Figure 5 --- the LLM-classical
crossover at super-rare is the headline F1 evidence.

\subsection{H4 --- Scaffolding \ensuremath{\times} case complexity (organ-system count)}\label{sec-7-8}

Pre-registered H4: multi-agent scaffolding \emph{helps on complex cases but hurts on
simple ones} (overthinking). Complexity = \# distinct HPO organ systems the gold
phenotype touches (single=1, oligo=2--3, multi=4+; HPO-input layers, full-N
2026-07-06). Scaffold \ensuremath{-} no-scaffold-control (llm\_control) R@1 delta, Gemini Flash:

\begin{table*}[t]
\centering\scriptsize\setlength{\tabcolsep}{3.5pt}
\caption{Multi-agent lift over the single-LLM control, by case complexity (H4).}\label{tbl:h4}
\begin{tabular}{@{}
lll@{}}
\toprule
Complexity & mdagents \ensuremath{-} control & medagents \ensuremath{-} control \\
\midrule
single-system (n\ensuremath{\approx}221--321) & \textbf{\ensuremath{-}0.08} & \textbf{\ensuremath{-}0.09} \\
oligo (2--3, n\ensuremath{\approx}765--1061) & \ensuremath{-}0.01 & \ensuremath{-}0.05 \\
multi-system (4+, n\ensuremath{\approx}3012--4329) & \textbf{+0.00} & \ensuremath{-}0.02 \\
\bottomrule
\end{tabular}
\end{table*}

\textbf{H4 supported (FWE-robust, \Cref{sec-8-8})}: the difference-of-differences ---
(mdagents\ensuremath{-}control on multi) \ensuremath{-} (mdagents\ensuremath{-}control on single) = +0.081 --- is
significant at full-N (2-prop z=2.51, Holm-adj p=0.012), up from the n=42
sub-bin pilot that was under-powered. The mechanism is a \emph{shrinking penalty}
rather than a gain: both multi-agent scaffolds clearly \emph{trail} the single-LLM
control on single-system cases (overthinking simple presentations, \ensuremath{-}0.08/\ensuremath{-}0.09)
and \emph{catch up to parity} (mdagents \ensuremath{\approx} +0.00) as organ-system involvement grows.
This sharpens \Cref{sec-7-2}: the small average scaffolding gain (F2) is really a \emph{penalty
on simple cases that dissolves on complex ones} --- the scaffolding's benefit is
avoiding its own overthinking cost when the case genuinely warrants deliberation,
not adding accuracy above the control.

\subsection{A6 --- Data-Contamination Audit (TS-Guessing approximation)}\label{sec-7-10}

Anticipated objection \#1 is that LLM agents may perform well only because
pre-cutoff PubMed corpora \emph{contain} the diseases we test on; the agents
would be exploiting training-frequency, not phenotype-disease reasoning.
We test this via a TS-Guessing approximation: for each gold ORPHA in
our phase4a predictions (top 600 by occurrence, \ensuremath{\geq}5 cases each), we
query NCBI PubMed \texttt{esearch.fcgi} with \texttt{"\textless{}disease\ name\textgreater{}"{[}All\ Fields{]}}
and \texttt{maxdate=2024/06/30} (a conservative pre-cutoff for all four
backbones), then correlate log(mention count + 1) with per-disease
R@1, per backbone, via Spearman \ensuremath{\rho}.

\begin{table*}[t]
\centering\scriptsize\setlength{\tabcolsep}{3.5pt}
\caption{Contamination correlation between pre-cutoff literature frequency and R@1, by backbone (H3/A6).}\label{tbl:contam}
\begin{tabular}{@{}
llll@{}}
\toprule
Backbone & n diseases & Spearman \ensuremath{\rho} & Interpretation \\
\midrule
\texttt{gemini} (Gemini 3 Flash) & 244 & \textbf{0.365} & weak \\
\texttt{gpt-5} (GPT-5 minimal) & 87 & \textbf{0.354} & weak \\
\texttt{v4-flash} (DeepSeek V4-Flash) & 244 & \textbf{0.348} & weak \\
\texttt{v4-pro} (DeepSeek V4-Pro) & 179 & \textbf{0.294} & weak \\
\texttt{lirical} (classical Bayesian) & 26 & \ensuremath{-}0.155 & null (control) \\
\texttt{vc\_rdagent} (offline IC) & 26 & \ensuremath{-}0.059 & null (control) \\
\bottomrule
\end{tabular}
\end{table*}

\textbf{Dichotomy} is the clean finding: every LLM backbone clusters at
\ensuremath{\rho} \ensuremath{\approx} 0.29--0.37, every classical / offline baseline at \ensuremath{\rho} \ensuremath{\approx} 0. The
classical baselines do not consume text --- they cannot have been
``trained on'' PubMed --- so their null \ensuremath{\rho} acts as a methodological control, confirming our pipeline introduces no spurious correlation.

\textbf{Reading}.
1. There IS a measurable training-frequency bias in LLM agents; the
simplest contamination critique is \emph{partially} supported.
2. But \ensuremath{\rho} \ensuremath{\approx} 0.3 means pre-cutoff exposure explains \ensuremath{\approx} 9 \% of R@1 variance
(\ensuremath{\rho}\textsuperscript{2} \ensuremath{\approx} 0.09), leaving \ensuremath{\approx} 91 \% to phenotype reasoning + extraction
quality + scaffold design. The contamination signal is real but
\textbf{bounded}.
3. The L4 post-cutoff PMC-OA holdout (2024-01-01+, after every backbone's
training cutoff) provides the bias-free reference. We now report the difficulty-matched cutoff experiment (H3, \Cref{sec-7-10-1}): performance does
\emph{not} drop across the training cutoff, bounding contamination from a
second, independent angle.

\subsubsection{H3 --- Difficulty-matched pre- vs post-cutoff (contamination, controlled)}\label{sec-7-10-1}

A naive ``post-cutoff R@1'' is confounded by dataset difficulty. We remove that
confound by building a pre-cutoff PMC set with the identical pipeline ---
same source (PMC-OA rare-disease case reports), same MeSH query, same
Gemini-3-Flash extraction, same Orphanet mapping, same Opus-4.8 gold
verification --- changing only the publication window (2016--2020, inside every backbone's training window vs 2024+, after all cutoffs). On the
Opus-diagnosis-agreed clean-gold subset (pre 195/220, post 198/198), Gemini-Flash:

\begin{table*}[t]
\centering\scriptsize\setlength{\tabcolsep}{3.5pt}
\caption{Difficulty-matched pre- vs.~post-cutoff R@1 (H3).}\label{tbl:h3}
\begin{tabular}{@{}
>{\raggedright\arraybackslash}p{(\linewidth - 6\tabcolsep) * \real{0.2500}}
  >{\raggedright\arraybackslash}p{(\linewidth - 6\tabcolsep) * \real{0.2500}}
  >{\raggedright\arraybackslash}p{(\linewidth - 6\tabcolsep) * \real{0.2500}}
  >{\raggedright\arraybackslash}p{(\linewidth - 6\tabcolsep) * \real{0.2500}}@{}}
\toprule
\begin{minipage}[b]{\linewidth}\raggedright
Agent
\end{minipage} & \begin{minipage}[b]{\linewidth}\raggedright
pre-cutoff R@1
\end{minipage} & \begin{minipage}[b]{\linewidth}\raggedright
post-cutoff R@1
\end{minipage} & \begin{minipage}[b]{\linewidth}\raggedright
\ensuremath{\Delta} (post\ensuremath{-}pre)
\end{minipage} \\
\midrule
llm\_control & 0.559 & 0.616 & +0.057 \\
mdagents & 0.582 & 0.611 & +0.029 \\
medagents & 0.564 & 0.626 & +0.062 \\
\textbf{pooled (single-pass)} & \textbf{0.568} & \textbf{0.618} & \textbf{+0.049} (z=1.72) \\
\bottomrule
\end{tabular}
\end{table*}

\textbf{Post-cutoff R@1 is at least as high as pre-cutoff for every agent.} If
memorisation inflated LLM rare-disease performance, memorisable (pre-cutoff)
cases would score \emph{higher} than unmemorisable (post-cutoff) ones --- the opposite
of what we observe. Strong performance transfers to genuinely unseen 2024+ cases,
so memorisation is not the driver. The small post-cutoff \emph{advantage} is most
plausibly newer case reports being marginally clearer (more routine genetic
confirmation), not contamination. Together A6 (weak within-dataset frequency
\ensuremath{\rho}\ensuremath{\approx}0.3) and H3 (no drop across the cutoff) bound contamination to a small effect.

\textbf{Strengthens F1, does not weaken it}. The LLMs' small \ensuremath{\rho}-explained
advantage is concentrated on diseases LLMs have seen more often; the
classical baselines deliver consistent reasoning across the entire
prevalence spectrum. This is the same direction as F1 (classical \textgreater{}
LLM on the rarest tier, \Cref{sec-7-7} H1, +28 pp on super-rare) viewed from a
different axis.

Visualised in Figure 4.

\subsection{H7 --- Failures cluster by specialty (shared blind spots)}\label{sec-7-9}

Pre-registered H7: agents' weakest specialties \emph{correlate across agents}
(Spearman \ensuremath{\rho}\ensuremath{\geq}0.6), implying dataset/ontology gaps rather than agent-specific
weaknesses. Specialty = modal HPO organ system per case (23-category HPO axis).
Cross-agent rank correlation of per-specialty R@1 (full-N 2026-07-06, 18 specialties n\ensuremath{\geq}10 each, Gemini Flash):

\begin{table*}[t]
\centering\scriptsize\setlength{\tabcolsep}{3.5pt}
\caption{Faithfulness-vs-accuracy rank correlation (H10).}\label{tbl:h10}
\begin{tabular}{@{}
ll@{}}
\toprule
Pair & Spearman \ensuremath{\rho} \\
\midrule
llm\_control vs mdagents & \textbf{0.93} \\
llm\_control vs medagents & \textbf{0.96} \\
mdagents vs medagents & \textbf{0.92} \\
\bottomrule
\end{tabular}
\end{table*}

\textbf{H7 confirmed and FWE-robust} (\Cref{sec-8-8}): all \ensuremath{\rho}\ensuremath{\geq}0.92 (up from the
n=13-specialty pilot at \ensuremath{\rho}\ensuremath{\approx}0.73); the conservative \ensuremath{\rho}=0.92 gives Holm-adj
p=0.0016. Universally weak specialties: nervous system
(0.11--0.14), metabolism/homeostasis (0.10--0.16), digestive (0.09); universally
\textbf{strong}: cardiovascular (0.41--0.44), integument (0.44--0.56). The shared
ordering points to ontology/data-level difficulty, not scaffold-specific gaps.
Notably the classical baselines invert the nervous-system weakness (LIRICAL
0.35, VC-RDAgent 0.43 vs LLM \textasciitilde0.12) and lead on head/neck (0.52--0.54) --- another
facet of F1's classical advantage. Visualised in Figure 7.


\begin{figure*}[t]
\centering
\includegraphics[width=0.82\textwidth]{figures/fig4_a6_contamination_scatter.png}
\caption{A6 TS-Guessing contamination scatter (LLM vs.\ classical).}\label{fig:fig4_a6_contamination_scatter}
\end{figure*}

\begin{figure*}[t]
\centering
\includegraphics[width=0.82\textwidth]{figures/fig5_prevalence_h1.png}
\caption{H1 prevalence-stratified R@1.}\label{fig:fig5_prevalence_h1}
\end{figure*}

\begin{figure*}[t]
\centering
\includegraphics[width=0.82\textwidth]{figures/fig6_hpo_density_h8.png}
\caption{H8 phenotype-density inverted-U.}\label{fig:fig6_hpo_density_h8}
\end{figure*}

\begin{figure*}[t]
\centering
\includegraphics[width=0.82\textwidth]{figures/fig7_specialty_h7.png}
\caption{H7 cross-agent specialty blind spots.}\label{fig:fig7_specialty_h7}
\end{figure*}

